% Format tugas kuliah berbasis pola umum template "homework/assignment"
% (article + identitas mahasiswa + daftar soal), mirip gaya Overleaf
% Professional Assignment / course problem set.
% Soal diadaptasi dari sumber daring sesuai susunan materi ppt/Logika_Predikat.tex.
\documentclass[12pt,a4paper]{article}

\usepackage[utf8]{inputenc}
\usepackage[T1]{fontenc}
\usepackage[indonesian]{babel}

\usepackage{geometry}
\geometry{margin=2.5cm}

% Font TeX Live/MiKTeX skalabel (menghindari konflik ekspansi font microtype + CM bitmap)
\usepackage{lmodern}

\usepackage{amsmath,amssymb,amsthm}
\usepackage{enumitem}
\usepackage{tabularx}
\usepackage{microtype}

\usepackage{fancyhdr}
\usepackage{lastpage}

\setlength{\headheight}{14pt}
\addtolength{\topmargin}{-2pt}

\usepackage[hidelinks,breaklinks]{hyperref}

% --- Isian identitas (ubah sesuai kebutuhan) ---
\newcommand{\JudulTugas}{Tugas 2}
\newcommand{\KodeMK}{TIF-xxx}
\newcommand{\NamaMK}{Logika Matematika}
\newcommand{\NamaMahasiswa}{[Nama Lengkap Mahasiswa]}
\newcommand{\NIMMhs}{[NIM]}
\newcommand{\Kelas}{[Kelas / Kelompok]}
\newcommand{\Prodi}{Teknik Informatika}
\newcommand{\Fakultas}{Fakultas Teknologi Informasi}
\newcommand{\Institusi}{Universitas Bale Bandung}
\newcommand{\DosenPengampu}{[Nama Dosen Pengampu]}
\newcommand{\TanggalKumpul}{[DD Bulan YYYY]}

\pagestyle{fancy}
\fancyhf{}
\fancyhead[L]{\small\NamaMahasiswa\ --- \NIMMhs}
\fancyhead[R]{\small\JudulTugas\ --- \NamaMK}
\fancyfoot[C]{\small Halaman \thepage\ dari \pageref*{LastPage}}
\renewcommand{\headrulewidth}{0.4pt}
\renewcommand{\footrulewidth}{0pt}

\setlength{\parindent}{0pt}
\setlength{\parskip}{0.6em}

% Jarak vertikal antar lingkungan soal (di atas/di bawah setiap \begin{soal}...\end{soal})
\newtheoremstyle{soaljarak}%
  {1.25ex plus 0.2ex minus 0.1ex}% ruang di atas
  {1.75ex plus 0.35ex minus 0.1ex}% ruang di bawah
  {\normalfont}% isi
  {0pt}%
  {\bfseries}%
  {.}%
  {0.5em}%
  {}
\theoremstyle{soaljarak}
\newtheorem{soal}{Soal}

\begin{document}

\begin{center}
  {\Large\bfseries \JudulTugas}\\[0.35em]
  {\large \NamaMK\ (\KodeMK)}\\[1em]
\end{center}

\noindent
\begin{tabularx}{\linewidth}{@{}p{4.8cm} @{\hspace{0.35em}:\hspace{0.65em}} >{\raggedright\arraybackslash}X@{}}
  Nama Mahasiswa            & \NamaMahasiswa \\
  Nomor Induk Mahasiswa     & \NIMMhs \\
  Program Studi             & \Prodi \\
  Fakultas                  & \Fakultas \\
  Institusi                 & \Institusi \\
  Kelas / Kelompok          & \Kelas \\
  Dosen Pengampu            & \DosenPengampu \\
  Batas tanggal pengumpulan & \TanggalKumpul \\
\end{tabularx}

\vspace{1em}
\hrule
\vspace{1em}

\section*{Petunjuk Umum}
\begin{itemize}[leftmargin=*]
  \item Jawaban \textbf{dianjurkan dikerjakan dalam format \LaTeX}; penggunaan \textit{Microsoft Word} atau aplikasi pengolah kata sejenis \textbf{dibolehkan}.
  \item Setelah selesai, \textbf{simpan atau ekspor jawaban menjadi berkas PDF}.
  \item \textbf{Unggah berkas PDF} tersebut ke \textbf{Edlink} sesuai tenggat di kontrak kuliah/RPS.
  \item Tugas ini membahas \textbf{logika predikat} (kalkulus predikat).
        Gunakan predikat, konstanta, fungsi, variabel, serta kuantor \(\forall\) (untuk semua) dan \(\exists\) (ada/untuk suatu) secara konsisten.
        Penghubung logik: \(\neg\), \(\wedge\), \(\vee\), \(\rightarrow\), \(\leftrightarrow\).
  \item \textbf{Terjemahan bahasa Indonesia \(\leftrightarrow\) formula:}
        perhatikan kata kuantifikasi (``semua'', ``setiap'', ``ada'', ``tidak semua'', ``hanya'', ``barang siapa'') dan urutan kuantor bersarang.
        Definisikan predikat/fungsi/konstanta yang Anda pakai sebelum menuliskan formula.
  \item \textbf{Term, proposisi, dan kalimat:}
        bedakan objek (term) dari pernyataan relasi (proposisi/kalimat).
        Sebutkan alasan singkat berdasarkan definisi pada materi kuliah.
  \item \textbf{Variabel bebas dan terikat:}
        tandai setiap kemunculan variabel; sebutkan kuantor mana yang mengikatnya (kuantor terdekat).
        Kalimat tertutup tidak memiliki variabel bebas.
  \item \textbf{Interpretasi:}
        nyatakan domain \(D\), lalu berikan nilai untuk setiap simbol bebas (konstanta, variabel bebas, fungsi, predikat).
        Tuliskan arti ekspresi dalam bahasa Indonesia berdasarkan interpretasi tersebut.
  \item \textbf{Aturan semantik kuantor:}
        evaluasi \(\forall x\,A\) dan \(\exists x\,A\) dengan interpretasi yang diperluas \(\langle x \leftarrow d \rangle \circ I\).
        Tunjukkan langkah penalaran, bukan hanya jawaban \(T\) atau \(F\).
  \item \textbf{Validitas:}
        hanya untuk kalimat tertutup.
        Untuk membuktikan valid: tunjukkan bernilai benar pada setiap interpretasi.
        Untuk membuktikan tidak valid: berikan \textbf{satu} interpretasi yang membuat kalimat bernilai salah.
  \item Tulis jawaban rapi per nomor soal dan per butir; gambar atau diagram domain boleh digambar tangan asalkan terbaca.
\end{itemize}

\section*{Soal}

\begin{soal}[\normalfont Sumber: diadaptasi dari AIMA FOL Exercises, Ex.~11]
  Diberikan kosakata berikut.
  \begin{itemize}[leftmargin=*]
    \item \(\mathit{Occupation}(p,o)\): orang \(p\) memiliki pekerjaan \(o\).
    \item \(\mathit{Customer}(p_1,p_2)\): orang \(p_1\) adalah pelanggan orang \(p_2\).
    \item \(\mathit{Boss}(p_1,p_2)\): orang \(p_1\) adalah atasan orang \(p_2\).
    \item Konstanta pekerjaan: \(\mathit{Doctor}\), \(\mathit{Surgeon}\), \(\mathit{Lawyer}\), \(\mathit{Actor}\).
    \item Konstanta orang: \(\mathit{Emily}\), \(\mathit{Joe}\).
  \end{itemize}
  Tuliskan pernyataan-pernyataan berikut dalam logika predikat.
  \begin{enumerate}
    \item Emily adalah ahli bedah atau pengacara.
    \item Joe adalah aktor, tetapi ia juga memiliki pekerjaan lain.
    \item Semua ahli bedah adalah dokter.
    \item Joe tidak memiliki pengacara (bukan pelanggan dari pengacara mana pun).
    \item Emily memiliki atasan yang merupakan pengacara.
    \item Ada seorang pengacara yang semua pelanggannya adalah dokter.
    \item Setiap ahli bedah memiliki seorang pengacara.
  \end{enumerate}
\end{soal}

\begin{soal}[\normalfont Sumber: diadaptasi dari Saarland Semantic Theory, Exercise~1.1; Rosen Discrete Mathematics §1.4--1.5]
  Terjemahkan kalimat-kalimat berikut ke dalam logika predikat.
  Definisikan predikat yang Anda gunakan.
  Usahakan mempertahankan struktur kalimat sebaik mungkin.
  \begin{enumerate}
    \item Budi pandai tetapi malas.
    \item Meskipun Budi adalah mahasiswa yang pandai, ia tidak akan lulus kecuali ia belajar keras.
    \item Setiap mahasiswa yang belajar keras akan lulus.
    \item Jika ada orang yang ribut, maka semua orang merasa terganggu.
    \item Jika seorang mahasiswa belajar keras, maka ia akan lulus.
    \item Paus adalah mamalia.
    \item Anjing yang menggonggong tidak menggigit.
    \item Setiap bilangan genap adalah kelipatan dari \(2\) (domain: bilangan bulat; gunakan predikat \(E(x)\): ``\(x\) genap'').
    \item Setiap mahasiswa memiliki setidaknya satu teman yang berkencan dengan penggemar sepak bola.
    \item Jika seseorang adalah orang tua dan laki-laki, maka ia adalah ayah dari seorang anak.
  \end{enumerate}
\end{soal}

\begin{soal}[\normalfont Sumber: diadaptasi dari Saarland Semantic Theory, Exercise~1.2; AIMA FOL Exercises, Ex.~13]
  Ubahlah formula berikut menjadi pernyataan bahasa Indonesia yang alami (hindari menyebut nama variabel secara eksplisit).
  Gunakan kunci terjemahan:
  \(j =\) John, \(m =\) Mary,
  \(S(x)\): \(x\) adalah mahasiswa,
  \(R(x)\): \(x\) kaya,
  \(L(x,y)\): \(x\) berhubungan dengan \(y\),
  \(F(x)\): \(x\) adalah burung,
  \(G(x)\): \(x\) terbang,
  \(\mathit{SpeaksLanguage}(x,l)\): \(x\) berbicara bahasa \(l\),
  \(\mathit{Understands}(x,y)\): \(x\) memahami \(y\).
  \begin{enumerate}
    \item \(S(j) \vee R(m)\)
    \item \(\forall x\,(R(x) \vee S(x))\)
    \item \(\exists x\,(\neg S(x) \wedge \forall y\,(\neg(x=y) \rightarrow S(y)))\)
    \item \(\forall x\,(S(x) \rightarrow \exists y\,(R(y) \wedge L(x,y)))\)
    \item \(\exists x\,\neg(F(x) \rightarrow (F(x) \wedge G(x)))\)
    \item
          \[
          \begin{aligned}
          &\forall x\,\forall y\,\forall l\,\bigl(
            \mathit{SpeaksLanguage}(x,l) \wedge \mathit{SpeaksLanguage}(y,l)\\
          &\qquad\rightarrow \mathit{Understands}(x,y) \wedge \mathit{Understands}(y,x)
          \bigr)
          \end{aligned}
          \]
  \end{enumerate}
  Selain itu, terjemahkan ke logika predikat (definisikan predikat yang dipakai):
  \begin{enumerate}[resume]
    \item Pemahaman menghasilkan persahabatan.
    \item Persahabatan bersifat transitif.
  \end{enumerate}
\end{soal}

\begin{soal}[\normalfont Sumber: diadaptasi dari AIMA FOL Exercises, Ex.~10]
  Pada setiap butir, diberikan kalimat bahasa Indonesia beserta beberapa kandidat ekspresi logika.
  Untuk setiap kandidat, nyatakan apakah ekspresi tersebut:
  (1)~benar-benar menyatakan makna kalimat;
  (2)~tidak valid secara sintaksis sehingga tidak bermakna; atau
  (3)~valid secara sintaksis tetapi tidak menyatakan makna kalimat.
  Predikat yang dipakai: \(\mathit{In}(x,y)\), \(\mathit{Borders}(x,y)\), \(\mathit{Country}(x)\), serta konstanta geografis yang sesuai.
  \begin{enumerate}
    \item Paris dan Marseille keduanya berada di Prancis.
    \begin{enumerate}
      \item \(\mathit{In}(\mathit{Paris} \wedge \mathit{Marseille}, \mathit{France})\)
      \item \(\mathit{In}(\mathit{Paris},\mathit{France}) \wedge \mathit{In}(\mathit{Marseille},\mathit{France})\)
      \item \(\mathit{In}(\mathit{Paris},\mathit{France}) \vee \mathit{In}(\mathit{Marseille},\mathit{France})\)
    \end{enumerate}
    \item Ada negara yang berbatasan dengan Irak dan Pakistan.
    \begin{enumerate}
      \item \(\exists c\,\bigl(\mathit{Country}(c) \wedge \mathit{Borders}(c,\mathit{Iraq}) \wedge \mathit{Borders}(c,\mathit{Pakistan})\bigr)\)
      \item \(\exists c\,\bigl(\mathit{Country}(c) \rightarrow (\mathit{Borders}(c,\mathit{Iraq}) \wedge \mathit{Borders}(c,\mathit{Pakistan}))\bigr)\)
      \item \(\bigl(\exists c\,\mathit{Country}(c)\bigr) \rightarrow \bigl(\mathit{Borders}(c,\mathit{Iraq}) \wedge \mathit{Borders}(c,\mathit{Pakistan})\bigr)\)
      \item \(\exists c\,\mathit{Borders}(\mathit{Country}(c), \mathit{Iraq} \wedge \mathit{Pakistan})\)
    \end{enumerate}
    \item Semua negara yang berbatasan dengan Ekuador berada di Amerika Selatan.
    \begin{enumerate}
      \item \(\forall c\,\bigl(\mathit{Country}(c) \wedge \mathit{Borders}(c,\mathit{Ecuador}) \rightarrow \mathit{In}(c,\mathit{SouthAmerica})\bigr)\)
      \item \(\forall c\,\bigl(\mathit{Country}(c) \rightarrow (\mathit{Borders}(c,\mathit{Ecuador}) \rightarrow \mathit{In}(c,\mathit{SouthAmerica}))\bigr)\)
      \item \(\forall c\,\bigl((\mathit{Country}(c) \rightarrow \mathit{Borders}(c,\mathit{Ecuador})) \rightarrow \mathit{In}(c,\mathit{SouthAmerica})\bigr)\)
      \item \(\forall c\,\bigl(\mathit{Country}(c) \wedge \mathit{Borders}(c,\mathit{Ecuador}) \wedge \mathit{In}(c,\mathit{SouthAmerica})\bigr)\)
    \end{enumerate}
    \item Tidak ada wilayah di Amerika Selatan yang berbatasan dengan wilayah mana pun di Eropa.
    \begin{enumerate}
      \item \(\neg\bigl(\exists c\,\exists d\,(\mathit{In}(c,\mathit{SouthAmerica}) \wedge \mathit{In}(d,\mathit{Europe}) \wedge \mathit{Borders}(c,d))\bigr)\)
      \item \(\forall c\,\forall d\,\bigl((\mathit{In}(c,\mathit{SouthAmerica}) \wedge \mathit{In}(d,\mathit{Europe})) \rightarrow \neg\mathit{Borders}(c,d)\bigr)\)
      \item \(\neg\forall c\,\bigl(\mathit{In}(c,\mathit{SouthAmerica}) \rightarrow \exists d\,(\mathit{In}(d,\mathit{Europe}) \wedge \neg\mathit{Borders}(c,d))\bigr)\)
      \item \(\forall c\,\bigl(\mathit{In}(c,\mathit{SouthAmerica}) \rightarrow \forall d\,(\mathit{In}(d,\mathit{Europe}) \rightarrow \neg\mathit{Borders}(c,d))\bigr)\)
    \end{enumerate}
  \end{enumerate}
\end{soal}

\begin{soal}[\normalfont Sumber: diadaptasi dari AIMA FOL Exercises, Ex.~12]
  Sama seperti Soal~4: klasifikasikan setiap kandidat sebagai (1), (2), atau (3).
  \begin{enumerate}
    \item Setiap kucing mencintai ibu atau ayahnya.
    \begin{enumerate}
      \item \(\forall x\,\bigl(\mathit{Cat}(x) \rightarrow \mathit{Loves}(x, \mathit{Mother}(x) \vee \mathit{Father}(x))\bigr)\)
      \item \(\forall x\,\bigl(\neg\mathit{Cat}(x) \vee \mathit{Loves}(x,\mathit{Mother}(x)) \vee \mathit{Loves}(x,\mathit{Father}(x))\bigr)\)
      \item \(\forall x\,\bigl(\mathit{Cat}(x) \wedge (\mathit{Loves}(x,\mathit{Mother}(x)) \vee \mathit{Loves}(x,\mathit{Father}(x)))\bigr)\)
    \end{enumerate}
    \item Setiap anjing yang mencintai salah satu saudaranya adalah bahagia.
    \begin{enumerate}
      \item \(\forall x\,\bigl(\mathit{Dog}(x) \wedge (\exists y\,(\mathit{Brother}(y,x) \wedge \mathit{Loves}(x,y))) \rightarrow \mathit{Happy}(x)\bigr)\)
      \item \(\forall x\,\forall y\,\bigl(\mathit{Dog}(x) \wedge \mathit{Brother}(y,x) \wedge \mathit{Loves}(x,y) \rightarrow \mathit{Happy}(x)\bigr)\)
      \item \(\forall x\,\bigl(\mathit{Dog}(x) \wedge (\forall y\,(\mathit{Brother}(y,x) \leftrightarrow \mathit{Loves}(x,y))) \rightarrow \mathit{Happy}(x)\bigr)\)
    \end{enumerate}
    \item Tidak ada anjing yang menggigit anak dari pemiliknya.
    \begin{enumerate}
      \item \(\forall x\,\bigl(\mathit{Dog}(x) \rightarrow \neg\mathit{Bites}(x, \mathit{Child}(\mathit{Owner}(x)))\bigr)\)
      \item \(\neg\exists x\,\exists y\,\bigl(\mathit{Dog}(x) \wedge \mathit{Child}(y,\mathit{Owner}(x)) \wedge \mathit{Bites}(x,y)\bigr)\)
      \item \(\forall x\,\bigl(\mathit{Dog}(x) \rightarrow \forall y\,(\mathit{Child}(y,\mathit{Owner}(x)) \rightarrow \neg\mathit{Bites}(x,y))\bigr)\)
      \item \(\neg\exists x\,\bigl(\mathit{Dog}(x) \rightarrow (\exists y\,(\mathit{Child}(y,\mathit{Owner}(x)) \wedge \mathit{Bites}(x,y)))\bigr)\)
    \end{enumerate}
  \end{enumerate}
\end{soal}

\begin{soal}[\normalfont Sumber: diadaptasi dari materi sintaks logika predikat (CS Waterloo / pola definisi pada slide kuliah)]
  Berdasarkan definisi simbol, term, proposisi, dan kalimat pada materi \textit{Logika Predikat}, tentukan untuk setiap ekspresi berikut apakah ia merupakan \textbf{term}, \textbf{proposisi}, \textbf{kalimat}, atau \textbf{bukan ekspresi yang terbentuk dengan baik}.
  Berikan alasan singkat.
  Asumsikan: \(a,b\) konstanta; \(x,y\) variabel; \(f\) fungsi unary; \(g\) fungsi biner; \(p\) predikat biner; \(q\) predikat unary.
  \begin{enumerate}
    \item \(f(a)\)
    \item \(p(a,x)\)
    \item \(g(x,f(a))\)
    \item \(\forall x\,q(x)\)
    \item \(p(f(a), g(x,y))\)
    \item \(\mathit{if}\; p(a,b)\; \mathit{then}\; f(a)\; \mathit{else}\; b\)
    \item \(\mathit{if}\; (\forall x)\,p(a,x)\; \mathit{then}\; (\exists y)\,q(y)\; \mathit{else}\; \neg p(a,b)\)
    \item \(p(a) \wedge f(x)\)
    \item \(\exists x\,p(x, f(x))\)
    \item \(q(p(a,b))\)
  \end{enumerate}
\end{soal}

\begin{soal}[\normalfont Sumber: diadaptasi dari pola latihan subekspresi pada slide \textit{Logika Predikat}]
  Diberikan ekspresi
  \[
  E :\;
  \mathit{if}\; (\forall x)\,q(x, f(a))\;
  \mathit{then}\; f(a)\;
  \mathit{else}\; b.
  \]
  \begin{enumerate}
    \item Sebutkan semua \textbf{subterm} dari \(E\).
    \item Sebutkan semua \textbf{subkalimat} dari \(E\).
    \item Sebutkan semua \textbf{subekspresi} dari \(E\).
  \end{enumerate}
  Selanjutnya, untuk ekspresi
  \[
  F :\;
  (\forall x)\,\bigl(p(x,y) \wedge (\exists y)\,r(y,z)\bigr),
  \]
  sebutkan semua subterm dan subkalimat dari \(F\).
\end{soal}

\begin{soal}[\normalfont Sumber: diadaptasi dari TUG Script Predicate Logic; CS Waterloo Predicate Logic Syntax]
  Untuk setiap formula berikut, tentukan setiap kemunculan variabel apakah \textbf{bebas} atau \textbf{terikat}.
  Jika terikat, sebutkan kuantor mana yang mengikatnya.
  Nyatakan pula apakah formula tersebut merupakan \textbf{kalimat tertutup}.
  \begin{enumerate}
    \item \((\forall x)\,[p(x,y) \wedge (\exists y)\,q(y,z)]\)
    \item \((\forall x)\,[p(x) \vee (\exists x)\,(\forall y)\,r(x,y)]\)
    \item \((\forall x)\,(\exists y)\,p(x,y)\)
    \item \((\forall x)\,p(x,y)\)
    \item \(\forall x\,\exists y\,P(x,y) \wedge Q(y,x)\)
      (perhatikan cakupan kuantor menurut aturan yang dipakai pada materi; jelaskan asumsi Anda)
    \item \((\forall x)\,(\exists y)\,(x \leq y) \vee (x \leq z)\)
  \end{enumerate}
\end{soal}

\begin{soal}[\normalfont Sumber: diadaptasi dari pola latihan interpretasi pada slide \textit{Logika Predikat}; AIMA FOL Exercises, Ex.~2--3]
  Diberikan ekspresi
  \[
  E = \mathit{IF}\; p(x,f(x))\; \mathit{THEN}\; (\exists y)\,p(a,y).
  \]
  \begin{enumerate}
    \item Misalkan \(I\) adalah interpretasi untuk \(E\) dengan domain bilangan real, di mana
          \(a = \sqrt{2}\), \(x = \pi\),
          \(f\) adalah fungsi ``dibagi dua'' (\(f_I(d) = d/2\)),
          dan \(p\) adalah relasi ``lebih besar atau sama dengan'' (\(p_I(d_1,d_2) = (d_1 \geq d_2)\)).
          Tuliskan arti ekspresi \(E\) berdasarkan \(I\), lalu tentukan nilai kebenarannya.
    \item Misalkan \(J\) adalah interpretasi untuk \(E\) dengan domain semua orang, di mana
          \(a =\) Soeharto, \(x =\) Soekarno,
          \(f\) adalah fungsi ``ibu dari'',
          dan \(p\) adalah relasi ``anak dari''.
          Tuliskan arti ekspresi \(E\) berdasarkan \(J\).
          (Nilai kebenaran boleh dibahas secara hipotetis jika fakta dunia nyata tidak Anda ketahui; yang penting adalah rumusan arti.)
    \item Pertimbangkan basis pengetahuan yang hanya berisi \(P(a)\) dan \(P(b)\).
          Apakah basis itu mengentail \(\forall x\,P(x)\)?
          Jelaskan dalam kerangka model/interpretasi.
    \item Apakah kalimat \(\exists x\,(x=x)\) valid?
          Jelaskan.
  \end{enumerate}
\end{soal}

\begin{soal}[\normalfont Sumber: diadaptasi dari Saarland Semantic Theory, Exercise~1.3--1.4]
  Pertimbangkan dua formula:
  \[
  A = \forall x\,(\mathit{Student}(x) \rightarrow \neg\mathit{PayAttention}(x)),
  \qquad
  B = \neg\exists x\,(\mathit{Student}(x) \wedge \mathit{PayAttention}(x)).
  \]
  \begin{enumerate}
    \item Bandingkan kondisi kebenaran \(A\) dan \(B\).
          Apakah keduanya ekuivalen?
          Jelaskan.
    \item Berikan satu struktur model di mana \(A\) bernilai benar, dan satu struktur model di mana \(A\) bernilai salah.
          Nyatakan domain serta ekstensi predikat \(\mathit{Student}\) dan \(\mathit{PayAttention}\).
    \item Kalimat ``setiap mahasiswa lulus'' secara intuitif sering dianggap mengimplikasikan ``ada mahasiswa yang lulus'', tetapi secara logika predikat tidak demikian.
          \begin{enumerate}
            \item Terjemahkan kedua kalimat ke logika predikat.
            \item Berikan sebuah model di mana terjemahan ``setiap mahasiswa lulus'' benar, tetapi terjemahan ``ada mahasiswa yang lulus'' salah.
          \end{enumerate}
  \end{enumerate}
\end{soal}

\begin{soal}[\normalfont Sumber: diadaptasi dari pola aturan semantik kuantor pada slide \textit{Logika Predikat}]
  Misalkan \(A : (\exists x)\,p(x,y)\).
  Diberikan interpretasi \(I\) dengan domain bilangan bulat positif,
  \(y = 2\), dan \(p_I(d_1,d_2) = (d_1 < d_2)\).
  \begin{enumerate}
    \item Tunjukkan langkah demi langkah bahwa \(A\) bernilai benar berdasarkan \(I\),
          dengan memakai aturan FOR SOME dan interpretasi yang diperluas.
    \item Selanjutnya, tinjau kalimat tertutup
          \[
          B :\;
          \mathit{IF}\; (\forall x)\,(\exists y)\,p(x,y)\;
          \mathit{THEN}\; p(a,f(a)).
          \]
          Misalkan \(I'\) meliputi domain bilangan real positif dengan
          \(a = 1\), \(f_{I'}(d) = \sqrt{d}\), dan \(p_{I'}(d_1,d_2) = (d_1 \neq d_2)\).
          Tentukan nilai kebenaran \(B\) berdasarkan \(I'\) dengan menelusuri antiseden dan konsekuen.
    \item Misalkan \(Q(x,y)\) menyatakan ``\(x + y = 0\)'' dengan domain bilangan real.
          Bandingkan nilai kebenaran \(\exists y\,\forall x\,Q(x,y)\) dan \(\forall y\,\exists x\,Q(x,y)\).
          Jelaskan mengapa urutan kuantor penting.
  \end{enumerate}
\end{soal}

\begin{soal}[\normalfont Sumber: diadaptasi dari AIMA FOL Exercises, Ex.~7; pola pembuktian validitas pada slide]
  Manakah di antara kalimat tertutup berikut yang \textbf{valid}?
  Untuk yang valid, berikan sketsa pembuktian berdasarkan aturan semantik.
  Untuk yang tidak valid, berikan satu interpretasi yang membuat kalimat bernilai salah.
  \begin{enumerate}
    \item \(\bigl(\exists x\,(x=x)\bigr) \rightarrow \bigl(\forall y\,\exists z\,(y=z)\bigr)\)
    \item \(\forall x\,\bigl(P(x) \vee \neg P(x)\bigr)\)
    \item \(\forall x\,\bigl(\mathit{Smart}(x) \vee (x=x)\bigr)\)
    \item \(\bigl[\neg(\forall x)\,p(x)\bigr] \leftrightarrow \bigl[(\exists x)\,\neg p(x)\bigr]\)
    \item \(\mathit{IF}\; (\exists y)\,(\forall x)\,q(x,y)\; \mathit{THEN}\; (\forall x)\,(\exists y)\,q(x,y)\)
  \end{enumerate}
\end{soal}

\begin{soal}[\normalfont Sumber: diadaptasi dari AIMA FOL Exercises, Ex.~14]
  Benar atau salah?
  Jelaskan jawaban Anda.
  \begin{enumerate}
    \item \(\exists x\,(x = \mathit{Rumpelstiltskin})\) adalah kalimat valid (selalu benar) dalam logika predikat.
    \item Setiap kalimat berkuantor eksistensial dalam logika predikat benar pada sembarang model yang berisi tepat satu objek.
    \item \(\forall x\,\forall y\,(x=y)\) dapat dipenuhi (\textit{satisfiable}).
    \item Basis pengetahuan \(\{P(a), P(b)\}\) mengentail \(\forall x\,P(x)\).
  \end{enumerate}
\end{soal}

\end{document}
